% Part: first-order-logic
% Chapter: natural-deduction
% Section: soundness-identity

\documentclass[../../../include/open-logic-section]{subfiles}

\begin{document}

\olfileid{fol}{ntd}{sid}

\olsection{Correctheid met \usetoken{S}{identity}}

\begin{prop}
Natuurlijke deductie met regels voor $\eq$ is correct.
\end{prop}

\begin{proof}
Elke !!{formula} van de vorm $\eq[t][t]$ is geldig, want voor elke
!!{structure}~$\Struct M$ geldt $\Sat{M}{\eq[t][t]}$. (Merk op dat we
veronderstellen dat de term $t$ gesloten is, dat wil zeggen dat hij
geen variabelen bevat, zodat bedelingen niet ter zake doen.)

Stel dat de laatste inferentie in !!a{derivation} \Elim{\eq} is, dat
wil zeggen dat de !!{derivation} de volgende vorm heeft:
\begin{prooftree}
  \AxiomC{$\Gamma_1$}
  \RightLabel{$\delta_1$}
  \DeduceC{$\eq[t_1][t_2]$}
  \AxiomC{$\Gamma_2$}
  \RightLabel{$\delta_2$}
  \DeduceC{$!A(t_1)$}
  \RightLabel{\Elim{\eq}}
  \BinaryInfC{$!A(t_2)$}
\end{prooftree}
De premissen $\eq[t_1][t_2]$ en $!A(t_1)$ worden respectievelijk
!!{derive}d uit de !!{undischarged} aannamen~$\Gamma_1$ en
$\Gamma_2$. We willen aantonen dat $!A(t_2)$ volgt uit $\Gamma_1
\cup \Gamma_2$. Beschouw !!a{structure}~$\Struct{M}$ waarvoor
$\Sat{M}{\Gamma_1 \cup \Gamma_2}$. Volgens de inductiehypothese geldt
$\Sat{M}{!A(t_1)}$ en $\Sat{M}{\eq[t_1][t_2]}$. Bijgevolg is
$\Value{t_1}{M} = \Value{t_2}{M}$. Laat $s$ een willekeurige bedeling
zijn en stel $m = \Value{t_1}{M} = \Value{t_2}{M}$. Volgens
\olref[fol][syn][ext]{prop:ext-formulas} geldt
$\Sat{M}{!A(t_1)}[s]$ dan en slechts dan als
$\Sat{M}{!A(x)}[\Subst{s}{m}{x}]$, en dit geldt dan en slechts dan als
$\Sat{M}{!A(t_2)}[s]$. Omdat $\Sat{M}{!A(t_1)}$, geldt dus
$\Sat{M}{!A(t_2)}$.
\end{proof}

\end{document}
